% =============================================================
% IEEE iSPEC 2026 Conference Paper  -- VERSION 4
% Rules: IEEEtran, two-column, strict 6 pages
% =============================================================
\documentclass[conference]{IEEEtran}
\IEEEoverridecommandlockouts
\usepackage{cite}
\usepackage{amsmath,amssymb,amsfonts}
\usepackage{algorithmic}
\usepackage{graphicx}
\usepackage{textcomp}
\usepackage{xcolor}
\usepackage{booktabs}
\usepackage{multirow}
\usepackage{url}

\def\BibTeX{{\rm B\kern-.05em{\sc i\kern-.025em b}\kern-.08em
    T\kern-.1667em\lower.7ex\hbox{E}\kern-.125emX}}

\begin{document}

\title{A Two-Stage Deep Learning Pipeline for Automated\\
Cell-Level Extraction and Defect Triage in Photovoltaic Modules%
\thanks{This work was carried out at [University Name].}}

% ================================================================
%  AUTHORS
% ================================================================
\author{%
\IEEEauthorblockN{1\textsuperscript{st} \textbf{[Author 1 Name]}}
\IEEEauthorblockA{\textit{\textbf{[Department]}} \\
\textit{Indian Institute of Technology Mandi}\\
Mandi, Himachal Pradesh, India \\
\textbf{[email@iitmandi.ac.in]}}
\and
\IEEEauthorblockN{2\textsuperscript{nd} \textbf{[Author 2 Name]}}
\IEEEauthorblockA{\textit{\textbf{[Department]}} \\
\textit{Indian Institute of Technology Mandi}\\
Mandi, Himachal Pradesh, India \\
\textbf{[email@iitmandi.ac.in]}}
}

\maketitle

\begin{abstract}
Automated visual inspection of photovoltaic (PV) modules is critical for scalable quality control, yet most published pipelines assume pre-cropped single-cell inputs and do not address the module-to-cell extraction problem. We present a two-stage deep learning pipeline that operates directly on raw module-level electroluminescence (EL) images. Stage~1 adopts YOLOv8-OBB (oriented bounding boxes) after a systematic study of fifteen classical OpenCV-based approaches, none of which generalised robustly. All five YOLOv8-OBB model scales achieve mAP@0.5\,=\,0.9950 on the test set, with the nano variant at 222~FPS. Stage~2 benchmarks seven architectures---ConvNeXt-Tiny, DenseNet-121, MobileNet-V3-Large, ResNet-101, EfficientNet-B4, Swin-S, and ViT-B/16---via Optuna hyperparameter search for binary EL defect triage (OK vs.\ B-grade). ConvNeXt-Tiny achieves best accuracy (97.2\,\%, AUC\,=\,98.6\,\%); DenseNet-121 leads on AUC/PR-AUC (98.9\,\%/99.3\,\%); MobileNet-V3-Large offers the best efficiency at 3.0M parameters.
\end{abstract}

\begin{IEEEkeywords}
photovoltaic inspection, electroluminescence, YOLOv8-OBB, oriented bounding boxes, binary defect classification, ConvNeXt, Optuna, deep learning
\end{IEEEkeywords}

%% ============================================================
\section{Introduction}

Global PV capacity is projected to exceed 10~TW by~2030~\cite{iea2023}, placing unprecedented demands on automated quality assurance. Electroluminescence (EL) imaging reveals recombination-active sub-surface defects invisible to conventional optics and has become the dominant non-destructive characterisation technique for commercial PV modules. Manual EL review cannot scale to modern manufacturing throughput or drone-based fleet inspection rates.

Automated EL defect analysis has been extensively studied at the \textit{single-cell} level. The canonical ELPV dataset~\cite{deitsch2019} provides 2{,}624 labelled cell-crop images; ResNet and EfficientNet backbones~\cite{he2016,tan2019} achieve binary accuracy exceeding 90\,\% on it. In realistic deployment, however, raw \textit{module-level} images arrive with arbitrary orientation, partial occlusion, and perspective distortion---no robust plug-in extraction stage exists.

Three unresolved gaps motivate this work: (1)~classical module-to-cell extraction fails to generalise across formats, contrast levels, and orientations; (2)~existing benchmarks compare few architectures without principled hyperparameter search; (3)~binary triage is a necessary precursor to fine-grained pixel-level defect segmentation.

\textbf{Contributions:}
\begin{enumerate}
  \item Systematic ablation of \textbf{fifteen classical OpenCV iterations} with root-cause analysis of four deployment-scale failure modes.
  \item \textbf{YOLOv8-OBB cell extractor} (mAP@0.5\,=\,0.9950, 222~FPS) covering full-module, partial-module, and single-cell inputs in one model.
  \item \textbf{7-architecture Optuna benchmark} for binary EL triage identifying ConvNeXt-Tiny as best overall and DenseNet-121 as best ranker.
  \item Roadmap for \textbf{Stage~3 pixel-level segmentation} with principled 29$\to$4 superclass taxonomy.
\end{enumerate}

%% ============================================================
\section{Related Work}
\label{sec:related}

\textbf{EL datasets.} The ELPV dataset~\cite{deitsch2019} (2{,}624 grayscale crops, $300\times300$~px) is the canonical benchmark. Its pre-cropped format has driven the single-cell paradigm, leaving module-level extraction unaddressed.

\textbf{Classical cell extraction.} Hough-transform line detection~\cite{duda1972}, morphological grid extraction~\cite{soille2003}, and projection profiling~\cite{wong1982} all rely on image-specific hyperparameters that do not transfer across formats, contrast levels, or orientations.

\textbf{Detection.} YOLO-family detectors~\cite{redmon2016,jocher2023} dominate real-time industrial inspection. YOLOv8-OBB extends axis-aligned detection with a learned rotation angle~$\theta$, enabling tight fitting of rotated objects.

\textbf{Classification.} ResNet~\cite{he2016} and EfficientNet~\cite{tan2019} dominate EL defect classification. ConvNeXt~\cite{liu2022} and Swin Transformer~\cite{liu2021} achieve superior fine-grained recognition but have not been benchmarked on EL with Optuna tuning~\cite{akiba2019}.

%% ============================================================
\section{Dataset}
\label{sec:dataset}

\subsection{Stage 1: Module-Level EL Images}
2{,}212 full-module JPEG EL images spanning 144-cell ($12\times12$) and 208-cell ($8\times26$) module formats from multiple manufacturers (ARTS, SDLE, BMRK subsets). Module-level split: 1{,}769 train / 443 val (80/20), preventing cell-level data leakage.

\subsection{Stage 2: Cell-Level Triage Images}
ELPV dataset~\cite{deitsch2019} augmented with Stage~1-extracted crops. Labels binarised: \textit{OK} vs.\ \textit{B-grade} (any defect). Module-level group split prevents spatial correlation across folds. Table~\ref{tab:datasplit} reports statistics. The test set is 32.6\,\% OK / 67.4\,\% B-grade, reflecting realistic deployment imbalance.

\begin{table}[t]
  \caption{Stage~2 dataset split statistics}
  \label{tab:datasplit}
  \begin{center}
  \renewcommand{\arraystretch}{1.05}\small
  \begin{tabular}{lrrr}
    \toprule
    \textbf{Split} & \textbf{Total} & \textbf{OK} & \textbf{B-grade}\\
    \midrule
    Train & 6{,}575 & 4{,}428 & 2{,}147\\
    Val   & 1{,}408 &   954   &   454\\
    Test  & 1{,}409 &   949   &   460\\
    \bottomrule
  \end{tabular}
  \end{center}
\end{table}

\subsection{Preprocessing}
Single-channel 8-bit grayscale $\to$ median blur ($3\times3$, noise suppression) $\to$ CLAHE (clipLimit\,=\,2.0, tileGridSize\,=\,$8\times8$, EL normalisation) $\to$ \texttt{cv2.warpPerspective} (canonical upright crops) $\to$ zero-mean unit-variance normalisation.

%% ============================================================
\section{Methodology}
\label{sec:method}

\subsection{Stage 1: Cell Extraction}
\label{sec:stage1}

\subsubsection{Classical OpenCV Exploration (15 notebooks)}

Two parallel tracks were studied. The \textbf{CellExtractOpenCV} track (pp1--pp11, 11 notebooks) targeted partial/single-cell images; the \textbf{FullModuleOpenCV} track (fmv1--fmv4, 4 notebooks) targeted full-resolution module images. Table~\ref{tab:classical_all} summarises all 15 iterations and their primary failure modes.

\begin{table}[t]
  \caption{Classical exploration: 15-notebook failure summary}
  \label{tab:classical_all}
  \begin{center}
  \renewcommand{\arraystretch}{1.0}\scriptsize
  \begin{tabular}{p{0.55cm}p{0.35cm}p{1.65cm}p{3.55cm}}
    \toprule
    \textbf{NB} & \textbf{Trk} & \textbf{Technique} & \textbf{Key Failure}\\
    \midrule
    pp1  & C & CLAHE+profiles   & Exploratory; no extraction\\
    pp2  & C & Morph.+warp      & Single image; not generalised\\
    pp3  & C & Inv.+contour     & Conflates cell/defect regions\\
    pp4  & C & Batch extractor  & \textbf{5--10\,\% skip rate}\\
    pp5  & C & ---              & YOLO pivot notebook\\
    pp6  & C & Gaussian+Canny   & Fragmented contours\\
    pp7  & C & NL-Means+Hough   & Image-specific threshold\\
    pp8  & C & Multi-cell Hough & Requires a priori line count\\
    pp9  & C & Otsu+Canny       & Format-specific count limit\\
    pp10 & C & Otsu+poly.\ fit  & Fails on multi-cell images\\
    pp11 & C & Hough+DBSCAN     & Dataset-specific $\epsilon$\\
    \midrule
    fmv1 & F & HoughLinesP      & Near-duplicate lines; unusable\\
    fmv2 & F & Proj.\ peaks     & \textbf{144 cells OK}; format-specific\\
    fmv3 & F & Batch fmv2       & 208 vs.\ 182 cells, same format\\
    fmv4 & F & DBSCAN merging   & Needs manual count verification\\
    \bottomrule
    \multicolumn{4}{l}{\scriptsize C\,=\,CellExtract; F\,=\,FullModule track.}
  \end{tabular}
  \end{center}
\end{table}

Four root-cause failure modes emerge across all 15 notebooks: (1)~\textit{Hyperparameter brittleness}---CLAHE limits, Hough thresholds, DBSCAN~$\epsilon$ tuned per image; (2)~\textit{Module format sensitivity}---projection-peak counting yielded 208 vs.\ 182 cells for nominally identical modules (fmv3); (3)~\textit{Contrast variability}---low-EL regions cause busbar lines to fall below adaptive threshold; (4)~\textit{Perspective intolerance}---axis-aligned morphological kernels break at even $5$--$10^\circ$ tilt.

\subsubsection{YOLOv8-OBB: Unified Final Approach}
\label{sec:yolo}

\textbf{Why not axis-aligned YOLO?} Standard axis-aligned YOLO detectors must \textit{fully enclose} rotated objects. For a $10^\circ$-tilted $144\times144$~px PV cell, the enclosing rectangle grows to $\approx$\textbf{[TBD]}~px, wasting $\sim$\textbf{[TBD]}\,\% of box area on background---degrading Stage~2 crop quality and causing heavy inter-cell NMS suppression. YOLOv8-OBB directly regresses $\theta$, keeping boxes tight to the cell boundary. YOLOv11-Seg is not compared: it is a segmentation model requiring pixel-level instance masks unavailable for Stage~1, and instance segmentation is the Stage~3 future-work target.

YOLOv8-OBB (Ultralytics~8.3.235~\cite{jocher2023}) was adopted as the sole Stage~1 detector. Each detection outputs a five-parameter oriented bounding box:
\begin{equation}
  \mathbf{d} = (c_x,\; c_y,\; w,\; h,\; \theta),
  \label{eq:obb}
\end{equation}
where $c_x, c_y \in \mathbb{R}$ are the centre pixel coordinates; $w, h \in \mathbb{R}^{+}$ are the width and height along the rotated axes; and $\theta \in [-90^\circ, 90^\circ)$ is the clockwise rotation angle of the box's major axis relative to the horizontal.

Post-processing converts each OBB to a rectified cell crop:
\begin{equation}
  (c_x, c_y, w, h, \theta)
  \xrightarrow{\;\text{corners}\;}
  M_{\text{persp}}
  \xrightarrow{\;\texttt{warpPersp.}\;}
  \hat{I}_{\text{cell}},
  \label{eq:postproc}
\end{equation}
where the \textit{corners} step maps OBB parameters to four corner coordinates $\{p_i\}_{i=1}^{4}$; $M_{\text{persp}} \in \mathbb{R}^{3\times3}$ is the perspective transformation matrix from \texttt{cv2.getPerspectiveTransform}; and $\hat{I}_{\text{cell}} \in \mathbb{R}^{H_0 \times W_0}$ is the canonical upright crop at fixed resolution $H_0 \times W_0$.

\subsection{Stage 2: Binary Quality Triage}
\label{sec:stage2}

\subsubsection{Benchmark Design}
Seven pretrained backbone families were fine-tuned end-to-end with a global average pooling head, dropout-regularised FC layer, and sigmoid binary output. No CBAM augmentation is applied in this baseline benchmark; attention-augmented variants are the immediate next step (Section~\ref{sec:conclusion}).

\subsubsection{Hyperparameter Search via Optuna}
Each backbone receives an identical Optuna trial budget; the search space is defined in Table~\ref{tab:optuna_space}. The objective is validation AUC (threshold-invariant under class imbalance). Fig.~\ref{fig:optuna_importance} shows the hyperparameter importance for ConvNeXt-Tiny; learning rate and weight decay dominate across all backbones.

\begin{table}[t]
  \caption{Optuna search space (Stage~2)}
  \label{tab:optuna_space}
  \begin{center}
  \renewcommand{\arraystretch}{1.05}\small
  \begin{tabular}{p{2.1cm}p{5.1cm}}
    \toprule
    \textbf{Hyperparameter} & \textbf{Search Space}\\
    \midrule
    Learning rate   & LogUniform$[10^{-5},\;10^{-2}]$\\
    Weight decay    & LogUniform$[10^{-6},\;10^{-2}]$\\
    Optimizer       & \{Adam, AdamW, SGD+mom.\}\\
    LR scheduler    & \{cosine, step, plateau\}\\
    Dropout         & Uniform$[0.01,\;0.5]$\\
    Batch size      & \{16, 32, 64\}\\
    Freeze strategy & \{none, partial, full\}\\
    Objective       & Validation AUC\\
    \bottomrule
  \end{tabular}
  \end{center}
\end{table}

\begin{figure}[t]
  \centering
  \includegraphics[width=\columnwidth]{fig_optuna_importance.png}
  \caption{Optuna hyperparameter importance (ConvNeXt-Tiny, Stage~2). Learning rate and weight decay are the dominant factors, consistent across all seven backbones.}
  \label{fig:optuna_importance}
\end{figure}

%% ============================================================
\section{Experimental Setup}
\label{sec:setup}

\textbf{Hardware:} NVIDIA RTX PRO 4500 Blackwell GPU; AMD Ryzen Threadripper 7960X (48 threads); 192~GB RAM; Ubuntu~24.04.2 LTS (Linux~6.17, 64-bit).

\textbf{Software:} Python~3.x, PyTorch~\textbf{[TBD]}, Ultralytics~8.3.235, \texttt{timm}~\textbf{[TBD]}, Optuna~\textbf{[TBD]}, OpenCV~4.12.0.88, NumPy~1.26.4, scikit-learn~\textbf{[TBD]}, Matplotlib~3.8.4.

\textbf{Reproducibility:} All experiments use a fixed random seed; results represent single-run performance. Multi-seed variance analysis is deferred to future work.

%% ============================================================
\section{Results}
\label{sec:results}

\subsection{Stage 1: Cell Extraction}

Table~\ref{tab:stage1_results} benchmarks all five YOLOv8-OBB scales on the 443-image test set. All variants achieve mAP@0.5\,=\,0.9950 with perfect recall (1.0000), confirming task saturation at the lightest scale. YOLOv8n-OBB (3.1M parameters, 4.5~ms/image) is the recommended deployment variant.

\begin{table}[t]
  \caption{Stage~1 YOLOv8-OBB benchmark (test set, 443 images)}
  \label{tab:stage1_results}
  \begin{center}
  \renewcommand{\arraystretch}{1.05}\small
  \begin{tabular}{lrrrr}
    \toprule
    \textbf{Model} & \textbf{Params} & \textbf{FPS} & \textbf{mAP@.5} & \textbf{mAP@.5:.95}\\
    \midrule
    YOLOv8n-OBB & 3.1M  & \textbf{222} & 0.9950 & 0.9909\\
    YOLOv8s-OBB & 11.4M & 216          & 0.9950 & 0.9932\\
    YOLOv8m-OBB & 26.4M & 148          & 0.9950 & \textbf{0.9947}\\
    YOLOv8l-OBB & 44.5M & 118          & 0.9949 & 0.9919\\
    YOLOv8x-OBB & 69.5M & 86           & 0.9950 & 0.9921\\
    \midrule
    Classical (fmv4) & --- & --- & N/A & N/A\\
    \bottomrule
    \multicolumn{5}{l}{\footnotesize fmv4: $\sim$90\,\% completeness (est.); no perspective correction.}
  \end{tabular}
  \end{center}
\end{table}

Fig.~\ref{fig:yolo_detections} shows YOLOv8n-OBB predictions on EL module images from the validation set.

\begin{figure}[t]
  \centering
  \includegraphics[width=\columnwidth,height=4.5cm,keepaspectratio=false]{fig_yolo_detections.jpg}
  \caption{YOLOv8n-OBB validation-batch predictions on full-module EL images. Oriented bounding boxes tightly localise individual cells regardless of module tilt; \texttt{cv2.warpPerspective} produces canonical upright crops for Stage~2 inference.}
  \label{fig:yolo_detections}
\end{figure}

\subsection{Stage 2: Binary Quality Triage}

Tables~\ref{tab:stage2_results}--\ref{tab:stage2_perclass} present the 7-architecture Optuna benchmark on 1{,}409 test images. ConvNeXt-Tiny achieves the highest accuracy (97.2\,\%) and MCC (93.7\,\%); DenseNet-121 achieves the highest AUC (98.9\,\%) and PR-AUC (99.3\,\%). MobileNet-V3-Large (3.0M parameters) matches ResNet-101 (42.5M) on accuracy, offering the best efficiency trade-off. Transformer-based Swin-S and ViT-B/16 underperform all CNN-based models. Both ConvNeXt-Tiny and DenseNet-121 substantially outperform the SVM (\textbf{[TBD]}\,\%) and CNN (\textbf{[TBD]}\,\%) baselines on ELPV reported by Deitsch et al.~\cite{deitsch2019}.

\begin{table}[t]
  \caption{Stage~2 binary triage: 7-architecture Optuna benchmark (test set; best trial per backbone)}
  \label{tab:stage2_results}
  \begin{center}
  \renewcommand{\arraystretch}{1.05}\small
  \setlength{\tabcolsep}{3.0pt}
  \begin{tabular}{p{1.75cm}rp{0.6cm}p{0.65cm}p{0.65cm}p{0.72cm}p{0.65cm}}
    \toprule
    \textbf{Model} & \textbf{Par.} & \textbf{Acc.} & \textbf{F1-M} & \textbf{AUC} & \textbf{PR-AUC} & \textbf{MCC}\\
    \midrule
    \textbf{ConvNeXt-T} & 27.8M & \textbf{.972} & \textbf{.968} & .986 & .989 & \textbf{.937}\\
    DenseNet-121        &  7.0M & .969           & .964          & \textbf{.989} & \textbf{.993} & .929\\
    MobileNet-V3L       &  3.0M & .960           & .954          & .970 & .973 & .909\\
    ResNet-101          & 42.5M & .960           & .953          & .980 & .988 & .908\\
    EffNet-B4           & 17.6M & .952           & .946          & .984 & .984 & .891\\
    Swin-S              & 48.8M & .949           & .940          & .976 & .986 & .884\\
    ViT-B/16            & 85.8M & .946           & .937          & .977 & .985 & .877\\
    \bottomrule
    \multicolumn{7}{l}{\footnotesize Par.\,=\,Parameters (M); F1-M\,=\,F1-Macro; MCC\,=\,Matthews Corr.\ Coeff.}
  \end{tabular}
  \end{center}
\end{table}

\begin{table}[t]
  \caption{Stage~2 per-class precision/recall (test set; OK\,=\,460, B-grade\,=\,949)}
  \label{tab:stage2_perclass}
  \begin{center}
  \renewcommand{\arraystretch}{1.05}\small
  \setlength{\tabcolsep}{3.0pt}
  \begin{tabular}{p{1.75cm}p{0.6cm}p{0.6cm}p{0.6cm}p{0.6cm}rr}
    \toprule
    \textbf{Model} & \textbf{OK-P} & \textbf{OK-R} & \textbf{BG-P} & \textbf{BG-R} & \textbf{FP} & \textbf{FN}\\
    \midrule
    \textbf{ConvNeXt-T} & .970 & .989 & .977 & .937 & 29 & 10\\
    DenseNet-121        & .974 & .980 & .958 & .946 & 25 & 19\\
    MobileNet-V3L       & .953 & .989 & .976 & .900 & 46 & 10\\
    ResNet-101          & .953 & .988 & .974 & .900 & 46 & 11\\
    EffNet-B4           & .958 & .972 & .940 & .913 & 40 & 27\\
    Swin-S              & .936 & .992 & .980 & .861 & 64 &  8\\
    ViT-B/16            & .936 & .987 & .971 & .861 & 64 & 12\\
    \bottomrule
    \multicolumn{7}{l}{\footnotesize FP\,=\,OK misclassified as BG; FN\,=\,BG misclassified as OK.}
  \end{tabular}
  \end{center}
\end{table}

Fig.~\ref{fig:stage2_roc} shows ROC curves; Fig.~\ref{fig:stage2_paramsf1} shows the parameters-vs.-F1 efficiency trade-off.

\begin{figure}[t]
  \centering
  \includegraphics[width=\columnwidth]{fig_stage2_roc.png}
  \caption{ROC curves (test set) for all seven Stage~2 architectures. DenseNet-121 (AUC\,=\,0.989) and ConvNeXt-Tiny (AUC\,=\,0.986) are the leading models; Transformer-based Swin-S and ViT-B/16 lag all CNN-based variants.}
  \label{fig:stage2_roc}
\end{figure}

\begin{figure}[t]
  \centering
  \includegraphics[width=\columnwidth]{fig_stage2_paramsf1.png}
  \caption{Parameters vs.\ test F1-macro for all seven Stage~2 architectures. MobileNet-V3-Large (3.0M, F1\,=\,0.954) achieves the best efficiency; ViT-B/16 (85.8M) provides the poorest return on model capacity.}
  \label{fig:stage2_paramsf1}
\end{figure}

%% ============================================================
\section{Discussion}
\label{sec:discussion}

\subsection{Stage 1: Task Saturation and Leakage Verification}

All five YOLOv8-OBB variants achieve mAP@0.5\,=\,0.9950. This saturation is consistent with PV cell grids: once the detector learns the periodic busbar-and-boundary signature from diverse training images, additional model capacity provides diminishing returns. Module-level partitioning (1{,}769 train / 443 val) explicitly prevents cell-level leakage---cells from the same module cannot appear in both splits. The nano variant's 4.5~ms/image latency makes real-time manufacturing-line integration feasible.

\subsection{Stage 2: CNN Inductive Bias on Small Data}

CNN-based models consistently outperform Transformers (Swin-S, ViT-B/16) across all metrics on the $\approx$6{,}575-sample training set. This is expected: CNNs' locality and translation-equivariance inductive biases provide stronger regularisation than self-attention on small data. ConvNeXt-Tiny and DenseNet-121 are statistically comparable on a single-seed run. ConvNeXt leads on accuracy and MCC, making it preferable for threshold-based production decisions; DenseNet leads on AUC and PR-AUC, making it preferable for ranking-based triage. The FN column in Table~\ref{tab:stage2_perclass} is critical for safety-sensitive applications: ConvNeXt minimises B-grade miss rate (FN\,=\,10) while Swin-S achieves FN\,=\,8 at the cost of 64 false alarms.

\textbf{End-to-end latency.} Stage~1 processes one module image in $\approx$4.5~ms (222~FPS). Stage~2 CNN forward pass on a $300\times300$ crop requires $<$\textbf{[TBD]}~ms on the same hardware. Full pipeline throughput is therefore dominated by the OBB detector and enables real-time deployment.

\subsection{Limitations}

\textit{Domain gap}: training on lab EL images; field-acquired EL introduces compression artefacts and ambient light. \textit{No CBAM yet}: attention-augmented variants are the immediate next experiment. \textit{Single-seed results}: multi-seed variance deferred. \textit{Annotation cost}: OBB labelling of 2{,}212 full-module images is labour-intensive; semi-supervised strategies are a natural follow-on.

%% ============================================================
\section{Conclusion and Future Work}
\label{sec:conclusion}

We presented a two-stage automated PV module inspection pipeline on raw module-level EL imagery. Fifteen classical OpenCV approaches exhibit four systematic failure modes motivating YOLOv8-OBB as a unified cell extractor; all five OBB scales achieve mAP@0.5\,=\,0.9950 at up to 222~FPS. A 7-architecture Optuna benchmark for binary triage identifies ConvNeXt-Tiny (Acc.\,97.2\,\%, AUC\,98.6\,\%) and DenseNet-121 (AUC\,98.9\,\%, PR-AUC\,99.3\,\%) as the leading architectures, with MobileNet-V3-Large as the best efficiency trade-off---all substantially exceeding ELPV baselines~\cite{deitsch2019}.

\textbf{Future work:}
(1)~\textit{CBAM augmentation}: channel and spatial attention after the final CNN stage, ablated against the baselines here.
(2)~\textit{Stage~3 pixel-level segmentation} on B-grade cells, consolidating 29 defect subtypes into four physically motivated superclasses (crack, inactive area, corrosion/metallisation, material/surface) and benchmarking U-Net variants against SegFormer~\cite{xie2021}.
(3)~Severity scoring calibrated against IV-curve power-loss measurements.
(4)~Domain adaptation from lab to field-acquired EL/IR imagery.
(5)~Embedded GPU benchmarking (Jetson Orin) for real-time manufacturing integration.

\bibliographystyle{IEEEtran}
\bibliography{references}

\end{document}